
\documentclass[
aps,
pra,
reprint,
superscriptaddress,
nofootinbib,
longbibliography
]{revtex4-2}

\usepackage{amsmath,amssymb}
\usepackage{bm}
\usepackage{hyperref}

\newcommand{\Hilb}{\mathcal H}
\newcommand{\Proj}{\operatorname{Proj}}
\newcommand{\Lat}{\mathcal L}

\begin{document}

\title{Contexts, Modalities, and Quantum Logic: A Critical Comparison of CSM with the Birkhoff--von Neumann Tradition}

\author{Karl Svozil}
\affiliation{Vienna, Austria}

\date{\today}

\begin{abstract}
The Contexts--Systems--Modalities (CSM) interpretation treats quantum
properties as definite modalities belonging to a system in a macroscopic
context.  This note compares that view with the quantum-logical method
introduced by Birkhoff and von Neumann and later sharpened by Gleason,
Specker, and Kochen.  Both approaches replace one global Boolean event algebra
by locally Boolean but mutually incompatible contexts.  The comparison is not
seamless, however.  Standard quantum logic is a type-I projection-lattice
formalism, whereas recent CSM emphasizes macroscopic, irreversible contexts
and type II/III operator algebras.  Moreover, the firefly logic $L_{12}$ shows
that contextual Boolean pasting alone is compatible with classical
two-valued states.  Thus CSM's primitives do not by themselves force
specifically quantum contextuality; the specifically quantum content enters
through the Hilbert-space or operator-algebraic representation, together with
Gleason--Uhlhorn and Kochen--Specker-type constraints.
\end{abstract}

\maketitle

\section{Introduction}

Grangier's recent note presents CSM as a non-FAPP, ontological interpretation
of quantum mechanics \cite{grangier2026}.  Its primitives are contexts,
systems, and modalities.  A modality is a definite, repeatable outcome, but it
is not a property of the isolated system alone; it belongs to a system together
with a definite macroscopic context.  The paper further stresses that real
contexts are not merely finite ancillary quantum systems: detectors and
records are macroscopic, dissipative, and irreversible, motivating an
operator-algebraic description of contexts.

This program has a natural point of contact with the quantum logic of
Birkhoff and von Neumann \cite{birkhoff1936}.  There, experimental yes--no
propositions are represented by closed subspaces, or equivalently by
projections, of a Hilbert space.  The resulting structure is an
orthocomplemented, orthomodular lattice and, in general, not a Boolean
algebra because distributivity fails.  Boolean logic survives only locally,
inside a maximal compatible context.  At this level the comparison is
immediate: CSM gives an ontological reading of a contextual structure that
quantum logic had already isolated algebraically.

The comparison nevertheless has to be made with some care.  CSM is not simply
``quantum logic plus ontology.''  The Birkhoff--von Neumann construction is
internal to a projection lattice such as $\Proj(\Hilb)$, normally the lattice
of a type-I Hilbert-space representation.  Recent CSM instead wants the
macroscopic context to be described by sectors, irreversible records, and
type II or type III von Neumann algebras \cite{haag1996,takesaki1979}.  Those
objects are not just larger Boolean blocks pasted inside the same type-I
lattice.  Thus the relation between CSM and quantum logic is both continuous
and tense: the vocabulary of contexts and modalities is naturally
quantum-logical, but the macroscopic algebraic setting changes the meaning of
some of the older lattice-theoretic notions.

\section{Quantum logic and CSM}

For a Hilbert space $\Hilb$, the quantum-logical event structure is
\begin{equation}
\Lat(\Hilb)=\Proj(\Hilb),
\end{equation}
partially ordered by inclusion of subspaces, with orthocomplement $P^\perp$.
Its nonclassical feature is not merely that distributivity can fail,
\begin{equation}
P\wedge(Q\vee R)\neq (P\wedge Q)\vee(P\wedge R),
\end{equation}
but that this failure occurs in the controlled setting of an orthomodular
lattice.  Orthomodularity is the residual order-theoretic structure left after
classical distributivity is abandoned; it is closely tied to projection,
compatibility, and the Sasaki-type conditional operations used in the
quantum-logical account of measurement \cite{kalmbach1983,ptak1991}.  More
concretely, when an ideal test of $Q$ gives the affirmative result, the
L\"uders rule
\begin{equation}
\rho\mapsto \frac{Q\rho Q}{\operatorname{Tr}(\rho Q)}
\end{equation}
is the state-space counterpart of this projection structure; in lattice
language the corresponding Sasaki projection sends $P$ to
$Q\wedge(P\vee Q^\perp)$.  Thus CSM's operational talk of changing context
meets quantum logic precisely at the point where projection, repeatability,
and conditionalization are tied together.  Within a maximal commuting family
of projections, however, the structure is Boolean.
Quantum logic is therefore not a rejection of ordinary logic inside a given
experiment; it is the rejection of one all-embracing Boolean algebra for all
possible experiments.  Specker famously illustrated this with the parable of
an overprotective seer who provides answers to individual questions but forbids
simultaneous answers to incompatible ones \cite{specker1960}.

This last point is often formulated more flexibly by using partial Boolean
algebras or pasted Boolean blocks.  In that language, one keeps Boolean
operations only inside compatible contexts and does not insist that all joins
and meets be globally meaningful.  Historically, this also brings the
comparison closer to Specker- and Kochen--Specker-style partial algebras than
to a single global orthomodular lattice, because CSM operationally needs only
context-wise Boolean structure plus cross-context identifications, not a fully
defined global meet and join for arbitrary propositions.  This is especially
close to the CSM picture, because a context in CSM also fixes a set of mutually
exclusive modalities, while changes of context need not be represented as
moves inside one classical event space.

Gleason's theorem supplies the probabilistic side of the type-I projection
structure: in Hilbert spaces of dimension at least three, every frame function
that is additive on each orthonormal resolution of the identity has the Born
form
\begin{equation}
\mu(P)=\operatorname{Tr}(\rho P).
\end{equation}
It is common to call this assumption probabilistic noncontextuality, since the
value assigned to a projector does not depend on the basis in which it is
embedded.  This use should be distinguished from Kochen--Specker
noncontextuality, which concerns $0$--$1$ truth-value assignments rather than
probability measures.  Specker's theorem and the Kochen--Specker theorem
provide the corresponding truth-value obstruction: suitable partial Boolean
or projection structures admit no global assignment of $0$'s and $1$'s
respecting the exclusivity constraints inside all contexts
\cite{specker1960,kochen1967}.  Specker's 1960 result is best read as a
partial-Boolean obstruction, whereas the Kochen--Specker theorem gives the
familiar Hilbert-space no-go theorem for noncontextual hidden variables.
The familiar two-dimensional exception should be kept in mind: neither the
original Gleason theorem nor Kochen--Specker obstruction applies directly to
qubits, so any fully general CSM reconstruction must either treat the qubit
case separately or appeal to enlarged measurement structures such as POVMs or
composite embeddings \cite{busch2003,caves2004}.

CSM translates this algebraic situation into ontological language.  In the
ordinary type-I Hilbert-space description, a context corresponds to a maximal
set of mutually orthogonal projectors resolving the identity.  A modality then
corresponds naturally to an atom, or definite proposition, in that Boolean
block.  A change of context is a change to another, generally incompatible,
Boolean block.  The additional CSM claim is that the context is not only a
basis in a calculation but a real macroscopic arrangement.  Hence the
definiteness of outcomes is not derived from unitary dynamics on one finite
Hilbert space; it is built into the ontology of system--context modalities
and, in Grangier's recent formulation, into the algebraic description of
irreversible macroscopic sectors \cite{grangier2026,haag1996}.

At exactly this point the comparison becomes delicate.  In type II and type III
factors there are no minimal projections; type III factors additionally admit
no nonzero normal semifinite trace. In particular, the normal-state structure
is very different from the atomic type-I case and does not support pure normal
states corresponding to lattice atoms.  Thus the lattice-theoretic phrase
``modality as an atom of a Boolean block'' cannot be carried unchanged into
the macroscopic algebraic regime.  If CSM identifies modalities in that regime
with sectors, folia, or equivalence classes of representations, then
``modality'' has shifted from an atomic proposition to an algebraic
superselection or representation-theoretic notion.  A sector, however, is not a
minimal yes--no event but a representation-theoretic class of states;
identifying it with a single definite outcome therefore requires an additional
interpretive step.  That shift may be physically well motivated, but it is not
part of the original Birkhoff--von Neumann lattice picture and should not be
hidden by vocabulary.

\begin{table*}
\caption{Compressed comparison of quantum logic and CSM.}
\label{tab:csm-ql}
\begin{tabular}{p{0.18\textwidth}p{0.25\textwidth}p{0.25\textwidth}p{0.24\textwidth}}
\hline\hline
Aspect & Quantum logic & CSM & Critical point \tabularnewline
\hline
Primitive & Propositions $P\in\Proj(\Hilb)$ & Modalities of system plus context & Same formal role in type I; different ontology \tabularnewline
Context & Boolean block or commuting set & Macroscopic arrangement & CSM makes the context physically real \tabularnewline
Classical domain & Boolean logic inside one block & Definite modalities inside one context & Essentially the same local Booleanity \tabularnewline
Nonclassical feature & No global Boolean algebra & No context-free list of properties & Largely the same obstruction, rephrased ontologically \tabularnewline
Probabilities & Gleason frame functions on projections & Transition weights between contexts & Formally linked, not competing accounts \tabularnewline
Truth values & Kochen--Specker obstruction & No global intrinsic modality assignment & Same no-go fact in lattice and ontic vocabularies \tabularnewline
Measurement update & Orthomodular/projection structure & Context change and repeatability & Natural meeting point, but rarely explicit in CSM \tabularnewline
Macroscopic side & Historically type-I focused; extendable to general von Neumann algebras & Irreversible sectors, type II/III algebras & Atom/modality translation becomes problematic \tabularnewline
Bell/EPR setting & Composite projection structure, but no apparatus ontology & Must treat separated macroscopic contexts & Relation between contextuality and nonlocality remains to be specified \tabularnewline
\hline\hline
\end{tabular}
\end{table*}

\section{Similarities, differences, and tensions}

The most important similarity is the replacement of global Booleanity by local
Booleanity.  Both approaches allow ordinary alternatives inside a single
context.  What they deny is the classical extrapolation from one Boolean
context to a universal Boolean algebra containing all possible alternatives.
For quantum logic this denial is expressed by non-distributivity,
orthomodularity, and the partial pasting of Boolean blocks.  For CSM it is
expressed by the thesis that modalities are meaningful only relative to a
specified system--context pair.

A second similarity is the role of probability.  In the
Birkhoff--von Neumann tradition, probabilities are states on the lattice of
propositions.  Gleason's theorem then shows that, under frame-function
additivity, such states are represented by density operators.  In CSM, the
same formal fact is read more physically: probabilities occur when a modality
prepared in one context is interrogated in another incompatible context.
Thus probabilities are not ignorance about a hidden Boolean state space, but
transition weights between contexts.  In the rank-one case these are literally
the usual Gleason--Born values: for modalities represented by
$P_\psi=|\psi\rangle\langle\psi|$ and
$P_\phi=|\phi\rangle\langle\phi|$,
\begin{equation}
\operatorname{Tr}(P_\psi P_\phi)=|\langle\phi|\psi\rangle|^2 ,
\end{equation}
which is exactly the transition weight CSM posits between the two modalities.
Grangier's program further claims that Born's rule can be recovered from
contextual quantization, using ingredients such as extracontextuality,
Uhlhorn-type symmetry considerations, and Gleason additivity
\cite{uhlhorn1963,auffeves2016,auffeves2017,grangier2022}.
This is also the point of Grangier's later discussion of Kolmogorovian
censorship and predictive incompleteness \cite{grangier2025kc}: within a
fixed context, or commuting subalgebra, the probabilities are ordinary
Kolmogorov probabilities, while the genuinely quantum question is how these
context-wise classical assignments are glued into one probability law on the
projection lattice.  In that language, the quantum state is predictively
incomplete until the measurement context is specified.  This reinforces the
quantum-logical reading of CSM: local classicality is unproblematic, and the
nonclassical content lies in the cross-context projection structure.
The same separation appears in finite non-Boolean event structures:
local normalization inside each context is weaker than single-valued
cross-context gluing, and even a glued admissible weight is still not thereby
classical, Born-rule realizable, or physically realized \cite{svozil2026softmax}.

This is also where a real limitation appears.  The derivation of Born's rule
does not follow from the bare words context, system, and modality.  It
requires a representation in which modalities can be identified across
contexts, transition probabilities have the required additivity or symmetry
properties, and the relevant Hilbert-space or operator-algebraic machinery is
already in place.  It is worth being explicit about the logical order.
Uhlhorn's theorem characterizes the symmetry maps of an \emph{already given}
orthogonality structure on rays, and Gleason's theorem presupposes the
projection lattice of a Hilbert space of dimension at least three.  Thus the
CSM derivation should be read as: \emph{given} that modalities are represented
by rays with the quantum orthogonality relation, the only admissible transition
weights are the Born ones.  The substantive quantum commitment is the ray
representation itself; once it is in place, Born's rule is close to forced.
This is precisely the asymmetry the firefly example exhibits at the qualitative
level.  In that sense the claimed ``derivation'' of Born's rule is conditional
rather than autonomous: once a Hilbert-space representation,
orthogonality-preserving context transformations, and Gleason-type additivity
are granted, the Born form follows; but those assumptions already encode much
of the specifically quantum structure.  CSM's distinctive ideas of
extracontextuality and extravalence are introduced precisely to bridge this
gap: they let the same modality be recognized across different contexts while
retaining contextual objectivity
\cite{auffeves2016,grangier2022,vandenbossche2023}.  From the
quantum-logical side, this is not merely an interpretive gloss; it is an
additional structural assumption.

The decisive difference is therefore not that quantum logic has mathematics
while CSM has ontology.  Rather, quantum logic is deliberately austere about
measurement apparatuses: it classifies propositions, compatibilities,
probability measures, and valuation obstructions, but it does not by itself
say what a macroscopic record is or why such a record is unique.  CSM makes
precisely this extra interpretive move.  It turns a context from a maximal
commuting family into a real macroscopic arrangement and treats a modality as
an actually possessed property of the system in that arrangement.

This difference explains the operator-algebraic emphasis in Grangier's recent
formulation.  The original quantum-logical scheme can be developed inside the
projection lattice of an ordinary Hilbert space, although quantum-logical
methods can also be generalized to projection lattices of von Neumann
algebras.  CSM claims that the type-I Hilbert-space picture is not enough for
the macroscopic side of measurement: real contexts are stable, irreversible,
and thermodynamic.  Hence type II or type III algebraic structures enter not
as mathematical decoration, but as a proposed description of the physical
definiteness of the apparatus.  Grangier's own illustration is instructive
here.  He couples each Stern--Gerlach output channel to an infinite reservoir
of two-level systems prepared in a KMS state, so that, in the intended
thermodynamic-limit idealization, the two post-measurement macrostates are
represented by inequivalent, ideally disjoint, GNS representations of the
quasi-local algebra \cite{grangier2026}.  This is a clean realization of the
claim that ``one result only'' is encoded by inequivalence rather than by
collapse.  But it also makes the lattice translation problem concrete: the
relevant ``modality'' is now a folium of states (a disjoint representation),
not an atom of a projection lattice.  For any finite reservoir, however, the
apparatus algebra remains type I and the relevant states lie in one
Hilbert-space representation; disjointness and strict irreversibility appear
only in the infinite-volume limit.

Finite-dimensional decoherence estimates make this contrast sharper rather
than removing it.  In a large environmental Hilbert space, typical branch
records can be quasi-orthogonal, with overlaps of order $e^{-S/2}$ when the
accessible dimension is $d_{\rm eff}\sim e^S$ \cite{svozil2026ortho}.  This
gives a powerful FAPP explanation of the invisibility of macroscopic
interference.  But without a specified system--environment split and
interaction Hamiltonian it does not uniquely select a pointer basis; and even
when einselection supplies stable pointer states, it does not turn an improper
mixture into a proper one.  It therefore supports the distinction drawn here:
ordinary Hilbert-space decoherence supplies record stability, while CSM's
ontology of macroscopic contexts is an additional interpretive commitment.

A natural question is how this contextual reading bears on Bell-type
nonlocality.  In the quantum-logical and partial-Boolean tradition, the
Kochen--Specker obstruction is state-independent and concerns a single system,
whereas Bell inequalities are state-dependent and concern spacelike-separated
wings.  CSM treats a Bell experiment as a single bipartite context: the joint
measurement arrangement at both wings fixes one set of mutually exclusive
modalities \cite{grangier2025kc}.  In that reading the ``no global valuation''
fact of Kochen--Specker and the ``no local hidden-variable'' fact of Bell are
closely related failures of one classical global probability space, though
Bell's theorem adds the specifically spacetime assumptions of local causality
and measurement independence.  What the firefly example shows for
contextuality has a Bell analogue: bare contextual vocabulary does not by
itself produce a Tsirelson-bounded violation; the quantitative violation again
comes from the Hilbert-space transition structure.

\section{\texorpdfstring{The firefly logic $L_{12}$}{The firefly logic L12}}

The firefly logic $L_{12}$ is a useful calibration point because it separates
weak contextuality from Kochen--Specker contextuality.  It consists of two
three-atom Boolean contexts,
\begin{equation}
C_1=\{a_1,a_2,a_3\},\qquad
C_2=\{a_3,a_4,a_5\},
\end{equation}
pasted along the common atom $a_3$.  Thus the Hasse diagram has twelve
elements and five atoms; operationally it is often described by a firefly in a
box observed from two sides, with the shared atom corresponding to the
``no flash'' alternative \cite{cohen1989,foulis1974,svozil1998,svozil2020}.

Quantum-logically, $L_{12}$ is already non-Boolean: the alternatives to the
shared atom depend on which test is performed.  In CSM language, this is a toy
model of two macroscopic contexts with one modality identifiable across both
contexts.  It therefore captures the grammar of contextual modalities in a
minimal way.

But $L_{12}$ is not a Kochen--Specker logic.  Because the two contexts share
only one atom, it admits many two-valued states, indeed enough to separate the
atoms.  Explicitly, either $v(a_3)=1$ and all other atoms receive value $0$, or
$v(a_3)=0$ and one chooses exactly one of $\{a_1,a_2\}$ and exactly one of
$\{a_4,a_5\}$ to receive value $1$. Thus $L_{12}$ has five dispersion-free
states, separating its five atoms.  One can locate $L_{12}$ on a finer scale.
Full Kochen--Specker sets forbid even a single global $\{0,1\}$ valuation
(state-independent contextuality), while weaker constructions only forbid
noncontextual reproduction of specific quantum statistics (state-dependent
contextuality).  $L_{12}$ fails both: it admits a full separating set of
two-valued states.  It is therefore not merely ``not state-independent KS''; it
is classically embeddable outright.  It is nondistributive and contextual as a
pasted empirical logic, yet still classically representable.  This is stronger
than the modest claim that contexts alone are insufficient for Kochen--Specker
contextuality.  It shows that the CSM primitives, taken by themselves, are
compatible with a classical hidden-variable representation.  The specifically
quantum content must be imported through the Hilbert-space or
operator-algebraic structure that CSM adds to those primitives.

The firefly example therefore marks a limit on the explanatory autonomy of
CSM.  A structure may have contexts, modalities, repeatability, and
non-distributive pasting while still admitting global two-valued states and no
Gleason-type Born rule.  Partition-logic analyses sharpen the same point:
the logico-empirical incidence structure can support classical probabilities
from the convex hull of dispersion-free states as well as Born-type
probabilities from faithful orthogonal graph representations, so the
structure alone does not select the probability theory
\cite{svozil2020graphs}.  To obtain the stronger CSM conclusion one needs
more: finite-dimensional Hilbert space with dimension at least three, or an
appropriate operator-algebraic analogue; extracontextual identification of
modalities across contexts; and the Gleason--Uhlhorn constraints on
transition probabilities.  The example is useful precisely because it keeps
these ingredients apart.

\section{Conclusion}

CSM can be read as a descendant of the Birkhoff--von Neumann insight, but not
as a smooth ontological completion of it.  Quantum logic already supplies much
of the formal grammar: local Boolean contexts, non-Boolean global structure,
orthomodularity, Gleason-type probability measures, and
Specker--Kochen restrictions on global truth assignments.  CSM accepts much
of this grammar but changes its interpretation: properties are modalities of
system--context pairs, and the context is a physically real macroscopic
arrangement.

The advantage of quantum logic is economy and precision.  It separates the
projection structure, probability theory, and valuation obstructions without
adding a detailed ontology of measurement.  The advantage of CSM is precisely
that it attempts to supply such an ontology, including the irreversibility and
unicity of macroscopic contexts.  But the attempt brings three obligations.
It must explain how type-I atomic modalities relate to type II/III sectors;
it must make clear which extra assumptions, beyond the bare CSM primitives,
recover Born's rule and exclude classical two-valued models; and it must say
how contextual objectivity reproduces Bell violations while specifying the
status of spacelike separated contexts and avoiding a hidden global context.

The firefly logic $L_{12}$ marks the lower boundary of this comparison.  It
displays contextual Boolean pasting, but not the full quantum obstruction.
Thus quantum logic gives the algebraic grammar of quantum propositions; CSM
proposes how that grammar is instantiated in actual measurements.  The
proposal is promising, but its force comes from the additional
Hilbert-space and operator-algebraic structure, not from contextual
vocabulary alone.  The real question is therefore not whether CSM is
contextual---so is quantum logic---but which additional structures turn
contextuality into specifically quantum ontology.

\begin{thebibliography}{99}

\bibitem{grangier2026}
P. Grangier,
``Revisiting the Interpretations of Quantum Mechanics:
From FAPP Solutions to Contextual Ontologies,''
arXiv:2601.20488v1 (2026).

\bibitem{birkhoff1936}
G. Birkhoff and J. von Neumann,
``The logic of quantum mechanics,''
Ann. Math. \textbf{37}, 823--843 (1936).

\bibitem{gleason1957}
A. M. Gleason,
``Measures on the closed subspaces of a Hilbert space,''
J. Math. Mech. \textbf{6}, 885--893 (1957).

\bibitem{specker1960}
E. Specker,
``Die Logik nicht gleichzeitig entscheidbarer Aussagen,''
Dialectica \textbf{14}, 239--246 (1960).

\bibitem{kochen1967}
S. Kochen and E. P. Specker,
``The problem of hidden variables in quantum mechanics,''
J. Math. Mech. \textbf{17}, 59--87 (1967).

\bibitem{kalmbach1983}
G. Kalmbach,
\emph{Orthomodular Lattices}
(Academic Press, London, 1983).

\bibitem{ptak1991}
P. Ptak and S. Pulmannova,
\emph{Orthomodular Structures as Quantum Logics}
(Kluwer, Dordrecht, 1991).

\bibitem{uhlhorn1963}
U. Uhlhorn,
``Representation of symmetry transformations in quantum mechanics,''
Ark. Fys. \textbf{23}, 307--340 (1963).

\bibitem{auffeves2016}
A. Auffeves and P. Grangier,
``Contexts, systems and modalities: a new ontology for quantum mechanics,''
Found. Phys. \textbf{46}, 121--137 (2016), arXiv:1409.2120,
doi:10.1007/s10701-015-9952-z.

\bibitem{auffeves2017}
A. Auffeves and P. Grangier,
``Recovering the quantum formalism from physically realist axioms,''
Sci. Rep. \textbf{7}, 43365 (2017), doi:10.1038/srep43365.

\bibitem{grangier2022}
P. Grangier,
``Revisiting quantum contextuality,''
arXiv:2201.00371 (2022).

\bibitem{grangier2025kc}
P. Grangier,
``Kolmogorovian Censorship, Predictive Incompleteness, and the locality loophole in Bell experiments,''
arXiv:2405.03184v3 (2025).

\bibitem{vandenbossche2023}
M. Van Den Bossche and P. Grangier,
``Revisiting quantum contextuality in an algebraic framework,''
J. Phys. Conf. Ser. \textbf{2533}, 012008 (2023), arXiv:2304.07757,
doi:10.1088/1742-6596/2533/1/012008.

\bibitem{busch2003}
P. Busch,
``Quantum states and generalized observables: a simple proof of Gleason's theorem,''
Phys. Rev. Lett. \textbf{91}, 120403 (2003),
doi:10.1103/PhysRevLett.91.120403.

\bibitem{caves2004}
C. M. Caves, C. A. Fuchs, K. K. Manne, and J. M. Renes,
``Gleason-type derivations of the quantum probability rule for generalized measurements,''
Found. Phys. \textbf{34}, 193--209 (2004),
doi:10.1023/B:FOOP.0000019581.00019.c1.

\bibitem{cohen1989}
D. W. Cohen,
\emph{An Introduction to Hilbert Space and Quantum Logic}
(Springer, New York, 1989).

\bibitem{foulis1974}
D. J. Foulis and C. H. Randall,
``Empirical logic and quantum mechanics,''
Synthese \textbf{29}, 81--111 (1974).

\bibitem{svozil1998}
K. Svozil,
\emph{Quantum Logic}
(Springer, Singapore, 1998).

\bibitem{svozil2020}
K. Svozil,
``Classical predictions for intertwined quantum observables are contingent and thus inconclusive,''
Quantum Rep. \textbf{2}, 278--292 (2020), arXiv:1707.08915.

\bibitem{svozil2020graphs}
K. Svozil,
``Faithful orthogonal representations of graphs from partition logics,''
Soft Comput. \textbf{24}, 10239--10245 (2020),
doi:10.1007/s00500-019-04425-1.

\bibitem{svozil2026ortho}
K. Svozil,
``The geometric part of decoherence:
Quasi-orthogonality in high-dimensional Hilbert spaces,''
in preparation (2026).

\bibitem{svozil2026softmax}
K. Svozil,
``Local softmax and global weights in non-Boolean event structures,''
in preparation (2026).

\bibitem{haag1996}
R. Haag,
\emph{Local Quantum Physics: Fields, Particles, Algebras},
2nd ed. (Springer, Berlin, 1996).

\bibitem{takesaki1979}
M. Takesaki,
\emph{Theory of Operator Algebras I}
(Springer, New York, 1979).

\end{thebibliography}

\end{document}
--- END OF FILE 3.tex ---

